\appendixsection[DG]{Auswertung Digitale Gestaltung (Tom Svilans)}
\label{anlage:interview-svilans}

Das Gespräch mit Tom Svilans konzentriert sich auf materialbasierte Entwurfsprozesse mit wiederverwendetem Holz: die digitale Erfassung gebrauchter Bauteile, den Umgang mit Unregelmäßigkeiten und den Erhalt ihrer architektonischen Qualitäten.

\subsection{Material als Ausgangspunkt des Entwurfs}

Wiederverwendetes Material bringt Spuren, Patina, frühere Verbindungen, Schäden, Maßabweichungen, Oberflächenqualitäten und konstruktive Einschränkungen in den Entwurfsprozess ein.

\begin{Blockquote}
The material already comes with a story, and that story starts to affect the design.
\end{Blockquote}

Das Werkzeug erfasst Material nicht nur über Länge, Breite, Höhe und Materialtyp. Es macht auch die gestalterischen, konstruktiven und konzeptionellen Eigenschaften eines Bauteils sichtbar.

Das digitale Werkzeug erfasst deshalb neben Geometrie und Materialtyp auch die Eigenschaften, die gestalterische, konstruktive und konzeptionelle Entscheidungen beeinflussen. Es unterstützt damit die Frage, welche neue Funktion ein konkretes Bauteil übernehmen kann.

\subsection{Vom Materiallesen zum digitalen Materialobjekt}

Der erste Schritt eines Reuse-Workflows ist das Lesen und Verstehen des vorhandenen Materials. Bevor ein Bauteil zugeordnet, bearbeitet oder eingebaut werden kann, muss geklärt werden, welche Eigenschaften es besitzt, welche Einschränkungen es mitbringt und welche Informationen fehlen.

\begin{Blockquote}
Before you design with it, you have to understand what the material can still do.
\end{Blockquote}

Diese frühe Phase der Materialbeobachtung bildet die Grundlage für die digitale Übersetzung. Wiederverwendete Bauteile lassen sich nicht ohne Weiteres wie neue Standardkomponenten in CAD-, BIM- oder parametrische Modelle einfügen. Sie sind oft verzogen, beschädigt, zugeschnitten, unvollständig dokumentiert oder nur teilweise bewertbar.

\begin{Blockquote}
The scan is not the answer; it only becomes useful when it helps you make a decision.
\end{Blockquote}

Das Reuse-Werkzeug benötigt deshalb eine Übersetzungsschicht zwischen physischem Material und digitalem Modell. Messungen, Fotos, Scans, Punktwolken, Meshes, Defekte, frühere Verbindungspunkte und Oberflächenzustände müssen in Informationen überführt werden, die im Entwurf tatsächlich verwendet werden können.

Das Ziel ist nicht zwingend ein perfekter digitaler Zwilling. Wichtiger ist ein brauchbares digitales Materialobjekt, das relevante Informationen bündelt: Geometrie und Orientierung, nutzbare und riskante Bereiche, frühere Verbindungen, Schäden, mögliche Anschlussbereiche, strukturelle Unsicherheiten sowie architektonisch relevante Oberflächen und Spuren.

\subsection{Unregelmäßigkeit und architektonische Qualität erhalten}

Ein wichtiger Punkt aus dem Interview ist, dass Unregelmäßigkeit nicht nur als Problem betrachtet werden sollte. Bei wiederverwendeten Materialien sind alte Ausschnitte, Verbindungen, Verformungen, Risse, Oberflächen und Maßabweichungen Teil der materiellen Realität. Sie können Einschränkungen erzeugen, aber auch Entwurfspotenziale eröffnen.

\begin{Blockquote}
You don't want to erase everything that makes the material interesting.
\end{Blockquote}

Diese Haltung betrifft nicht nur technische Bearbeitung, sondern auch architektonische Qualität. Wiederverwendung verändert die ästhetische, räumliche und narrative Qualität eines Projekts.

\begin{Blockquote}
Reuse is not only a technical problem; it changes the architectural character of the project.
\end{Blockquote}

Das digitale Werkzeug muss diese Perspektive abbilden. Viele Materialdatenbanken reduzieren Bauteile auf technische Parameter: Maße, Materialtyp, Zustand, Kosten, Gewicht oder Tragfähigkeit. Bei wiederverwendeten Materialien reicht das nicht aus. Gerade die nicht-standardisierten Eigenschaften können für den Entwurf entscheidend sein: Patina, Spuren, Farbe, frühere Fügungen, Gebrauchsspuren und sichtbare Geschichte.

Ein gutes Reuse-Werkzeug sollte deshalb Unterschiede nicht glätten, sondern lesbar machen. Es sollte sowohl technische Eigenschaften als auch architektonische Qualitäten aufnehmen können, damit entschieden werden kann, welche Spuren erhalten, sichtbar gemacht, überarbeitet oder bewusst in den Entwurf integriert werden.

\subsection{Entscheidungen unter Unsicherheit ermöglichen}

Wiederverwendete Materialien sind fast immer mit unvollständigem Wissen verbunden. Herkunft, Alter, frühere Belastung, Feuchtigkeit, Schäden, Kontamination, Tragfähigkeit oder zukünftige Nutzbarkeit sind oft nicht vollständig bekannt.

\begin{Blockquote}
You never know everything about reused material, but you still need to make decisions.
\end{Blockquote}

Das digitale Werkzeug macht Unsicherheit sichtbar. Das System sollte sichtbar machen, welche Informationen gesichert sind, welche auf Annahmen beruhen und wo weitere Prüfung notwendig ist.

Ein Bauteil könnte deshalb neben Geometrie und Materialtyp auch einen Prüfstatus, ein Vertrauensniveau, offene Risiken und empfohlene nächste Schritte enthalten. Ziel ist nicht, menschliche Beurteilung zu ersetzen, sondern Entscheidungen transparenter und nachvollziehbarer zu machen.

Damit unterscheidet sich ein entscheidungsunterstützendes Reuse-Werkzeug von einem einfachen digitalen Inventar: Es verwaltet nicht nur Objekte, sondern auch Wissensstände und offene Fragen.

\subsection{Materialzuordnung statt reiner Materialverwaltung}

Ein weiteres zentrales Thema aus dem Interview ist die Zuordnung vorhandener Materialien zu neuen Entwurfsaufgaben. Bei wiederverwendeten Bauteilen reicht es nicht, zu wissen, was vorhanden ist. Entscheidend ist, wofür ein vorhandenes Element geeignet sein könnte.

\begin{Blockquote}
The question is not only what material is available, but where it makes sense to use it.
\end{Blockquote}

Ein digitales Werkzeug müsste deshalb zwischen Materialinventar und Entwurfsentscheidung vermitteln. Es sollte Varianten vergleichbar machen: Welches Element passt zu welcher Anforderung? Welche Zuordnung reduziert Verschnitt? Wo ist ein Bauteil mit unklarer Qualität trotzdem sinnvoll einsetzbar? Welche Lösung erhält alte Oberflächen oder bestehende Verbindungen? Welche Option ist für Fertigung, Montage oder spätere Demontage am geeignetsten?

Das Werkzeug unterstützt die Materialzuordnung. Es unterstützt nicht nur die Verwaltung vorhandener Elemente, sondern hilft dabei, ihre mögliche Rolle im Entwurf zu bewerten.

\subsection{Wiederverwendung für den nächsten Lebenszyklus vorbereiten}

Ein wichtiger Gedanke aus dem Gespräch ist, dass Wiederverwendung nicht als einmalige Rettung eines Materials verstanden werden sollte. Wenn ein Bauteil in ein neues Projekt eingebaut wird, entsteht bereits die Frage nach seiner nächsten möglichen Nutzung.

\begin{Blockquote}
If we reuse something now, we should not make it harder to reuse it again later.
\end{Blockquote}

Der Materialdatensatz bleibt fortschreibbar. Er sollte nicht nur den aktuellen Zustand dokumentieren, sondern auch neue Bearbeitungsschritte, Schnitte, Verbindungen, Einbauorte, Prüfungen und Demontagemöglichkeiten aufnehmen.

Damit wird aus dem digitalen Materialobjekt eine Art Materialgedächtnis. Es unterstützt nicht nur die aktuelle Wiederverwendung, sondern bereitet zukünftige Wiederverwendung vor.

\subsection{Kritische Einordnung}

Die Ergebnisse des Interviews beziehen sich vor allem auf Holz und holzbasierte Arbeitsabläufe. Andere Materialgruppen — etwa Stahl, Beton, Glas, Fassadenelemente oder technische Komponenten — bringen andere Prüf-, Sicherheits-, Logistik- und Regulierungsfragen mit sich.

Trotzdem sind die methodischen Prinzipien übertragbar. Wiederverwendetes Material muss lesbar gemacht, digital übersetzt, bewertet, zugeordnet und für zukünftige Nutzung dokumentiert werden. Die konkreten Datenfelder und Prüfprozesse würden je nach Material variieren, aber die grundlegende Logik bleibt relevant.

Eine zweite kritische Frage betrifft die Gefahr der Überstandardisierung. Wenn digitale Werkzeuge wiederverwendete Materialien nur in saubere Datenobjekte übersetzen, könnten genau jene Qualitäten verloren gehen, die Reuse architektonisch interessant machen: Spuren, Patina, Unregelmäßigkeit, Geschichte und materielle Eigenheit.

Ein gutes Reuse-Werkzeug sollte deshalb nicht versuchen, gebrauchtes Material vollständig zu normalisieren. Es sollte helfen, Unterschiede sichtbar, interpretierbar und nutzbar zu machen.
